\chapter{Patterns: undo with branches}
\label{ch:undobranches}

\begin{aside}
Linear undo discards the future when you make a new change.
Branching undo keeps every branch as a tree. The user can
explore alternatives.
\end{aside}

\section{The tree}

\begin{lstlisting}
struct Node {
    Action action;
    Node* parent;
    std::vector<Node*> children;
};
Node* current;
\end{lstlisting}

The current node is the present. Going back walks toward
the root; going forward walks to a child.

\section{Branching}

A new action under \code{current} creates a new child.
\code{current} advances to it. Existing children remain
accessible.

\section{Navigation}

\begin{itemize}[leftmargin=*]
  \item Undo: go to parent.
  \item Redo: go to last-visited child.
  \item Branch picker: choose among siblings.
\end{itemize}

\section{Visualisation}

Vim has \code{:undolist}. \maya{} apps can show a tree
view of the history.

\section{Memory}

A long-lived editing session can accumulate a deep tree.
Compress: drop branches older than N steps; merge linear
runs.

\section{Recap}

\begin{recap}
\begin{itemize}[leftmargin=*]
  \item Tree of actions, not stack.
  \item Branches preserved on new action.
  \item Navigate parent/child/sibling.
  \item Compress old branches.
\end{itemize}
\end{recap}

\section*{Workshop}

\begin{enumerate}[leftmargin=*]
  \item Replace your undo stack with a tree.
  \item Add a branch picker.
  \item Implement compression.
\end{enumerate}
